\section{Conclusion}

\subsection{Summary of Contributions}

This study presents a controlled experimental \emph{framework} for evaluating performance–cost trade-offs in Amazon DynamoDB under serverless workloads, operationalised through a complete Infrastructure-as-Code artefact. Pipeline unit tests exercise moto; \textbf{no filled campaign cells} are claimed. The research makes three principal contributions:

\begin{enumerate}
    \item \textbf{Factorial experimental design:} A $3 \times 2 \times 4$ factorial structure simultaneously varying partition-key design (simple, composite, write-sharded), capacity mode (on-demand, provisioned), and workload profile (read-heavy, write-heavy, mixed, burst), with explicit measurement of throttling rate, consumed capacity, and cost per 10,000 operations—dependent variables absent from self-hosted benchmarks \citep{pantelic2026benchmarking}.
    
    \item \textbf{Reusable artefact:} Infrastructure-as-Code (Terraform), Lambda-based workload generators with Zipfian access distributions, CloudWatch metrics collectors, and statistical analysis scripts (two-way ANOVA with interaction, Tukey HSD post-hoc tests) are provided as open-source components, enabling third-party replication and extension to other DynamoDB configurations or cloud database services.
    
    \item \textbf{Methodological foundation for cloud-metered evaluation:} The study establishes a reproducible methodology for benchmarking cloud databases under realistic serverless workloads, addressing the gap between self-hosted benchmarks (which lack metered cost and admission control) and vendor-provided guidance (which lacks controlled empirical validation).
\end{enumerate}

\subsection{Anticipated Findings (Production Validation Required)}

This report does \textbf{not} present filled empirical campaign results. The following are design hypotheses for a future live AWS run:

\begin{enumerate}
    \item \textbf{Partition-key design effects:}
    \begin{itemize}
        \item K1 (simple keys) and K2 (composite keys) are expected to exhibit similar latency and cost under uniformly distributed access, but K2 enables range queries.
        \item K3 (write-sharded keys) should reduce write throttling under hot-key workloads but incur 10× read amplification, increasing read latency and cost.
    \end{itemize}
    
    \item \textbf{Capacity mode effects:}
    \begin{itemize}
        \item Provisioned mode is expected to offer lower cost under sustained workloads (W1–W3) but experience throttling during burst (W4) if auto-scaling response time lags traffic spikes.
        \item On-demand mode should eliminate throttling but incur 2–5× higher per-operation costs \citep{haubenschild2025oltp}.
    \end{itemize}
    
    \item \textbf{Interaction effects:}
    \begin{itemize}
        \item The latency-minimising configuration may differ from the cost-minimising configuration, requiring practitioners to make explicit trade-offs.
        \item Adaptive capacity \citep{aws2025adaptive} may enable provisioned-mode tables to absorb short bursts without throttling, narrowing the performance gap with on-demand mode.
    \end{itemize}
\end{enumerate}

These hypotheses remain to be tested on production infrastructure; simulation results do not constitute empirical evidence.

\subsection{Practical Implications (deferred)}

No practitioner configuration advice is asserted from this report's empty cells. After a live K1--K3 campaign fills latency/throttle/cost tables, architects may revisit trade-offs such as provisioned vs on-demand under burst and write-sharding read amplification. Until then, the lists below are \textbf{out of scope as evidence}:

\begin{itemize}
    \item Any claim that K1/K2/K3 or a capacity mode is empirically best for a named workload class.
    \item Any Cost Explorer--validated bill comparison (none committed).
    \item Any K4 adaptive-sharding ``dominance'' narrative (off-factorial; quarantined).
\end{itemize}

\subsection{Limitations}

This study has several limitations that constrain the generalisability of findings:

\begin{enumerate}
    \item \textbf{Single database service:} Evaluation focuses exclusively on Amazon DynamoDB; findings may not generalise to Azure Cosmos DB, Google Cloud Firestore, or self-hosted NoSQL databases due to differences in partitioning algorithms, billing models, and auto-scaling policies.
    
    \item \textbf{Synthetic workloads:} While Zipfian access distributions ($s = 1.070916$) model realistic skew observed in prior literature \citep{kanellis2025f2, ziegler2023oltp}, real-world applications may exhibit multi-modal distributions, temporal correlations, or access patterns not captured by a single Zipfian distribution.
    
    \item \textbf{Fixed item size:} All experiments use 1 KB items; applications with variable item sizes (1 byte to 400 KB) will experience different RCU/WCU consumption patterns and may encounter item-size-based throttling at partition limits (3,000 RCU/s, 1,000 WCU/s per partition).
    
    \item \textbf{Single region:} Evaluation is conducted in \texttt{eu-west-1}; multi-region replication (DynamoDB Global Tables) introduces additional cost ($\times$2 write cost) and latency (cross-region replication lag) not addressed here.
    
    \item \textbf{Eventually consistent reads:} All reads use eventual consistency; applications requiring strong consistency (\texttt{ConsistentRead=true}) consume 2× RCU per read, doubling read costs.
    
    \item \textbf{Empty empirical cells:} Core factorial tables remain placeholders; moto unit tests do not replace a live DynamoDB campaign. Production validation is required before applying any performance--cost guidance.
\end{enumerate}

\subsection{Future Work}

Several research directions extend this study:

\begin{enumerate}
    \item \textbf{Variable item sizes:} Extend the factorial design to include 1 KB, 8 KB, and 32 KB item sizes, testing whether write sharding (K3) remains beneficial when the hottest key approaches partition write limits (1,000 WCU/s $\times$ item-size-in-KB).
    
    \item \textbf{Multi-region replication:} Evaluate DynamoDB Global Tables under the same workload profiles, quantifying cross-region latency, replication lag, and cost multipliers for applications requiring geographic distribution.
    
    \item \textbf{Composite access patterns:} Investigate workloads with multiple Zipfian distributions (e.g., daily trending keys + long-tail historical keys) or temporal correlations (time-series data with hot recent keys).
    
    \item \textbf{Transactional workloads:} Extend the evaluation to DynamoDB transactions (\texttt{TransactWrite}, \texttt{TransactRead}), which consume 2× capacity and introduce cross-item coordination latency \citep{idziorek2023transactions}.
    
    \item \textbf{Cross-service comparison:} Replicate the factorial design on Azure Cosmos DB (with throughput units and partition-key design) and Google Cloud Firestore (with automatic sharding) to identify service-specific trade-offs.
    
    \item \textbf{Machine learning-driven sharding:} Integrate learned sharding policies (e.g., Neuroshard \citep{eldeeb2022neuroshard}) that dynamically adjust shard counts based on observed access patterns, comparing learned policies to static K3 write sharding.
    
    \item \textbf{Cost-performance optimisation:} Develop a multi-objective optimisation framework (Pareto frontier) that selects table configurations to minimise cost subject to latency SLAs, accounting for workload variability and item-size distributions.
\end{enumerate}

\subsection{Final Remarks}

Cloud-native databases such as Amazon DynamoDB have abstracted operational complexity (server provisioning, replication, scaling) but have introduced new configuration complexity: partition-key design and capacity mode decisions that fundamentally determine performance and cost. This study contributes a rigorous experimental \emph{framework} for evaluating these configuration-level trade-offs, with artefact plumbing covered by moto unit tests and designed for third-party replication on production infrastructure. Extending Panteli\'{c} et al.\ latency/throughput measures to metered dimensions (throttling, consumed capacity, cost) is the intended CA2 contribution; \textbf{filled live campaign cells are the sole remaining empirical residual}.

The complete artefact (Infrastructure-as-Code, workload generators, analysis scripts) is provided to enable reproducibility. Live DynamoDB execution under the alignment-first gate is the next step.
